\documentclass[11pt]{article}

\usepackage[margin=1in]{geometry}
\usepackage{amsmath,amssymb,amsthm}
\usepackage{mathtools}
\usepackage{hyperref}
\usepackage{xcolor}
\usepackage{enumitem}
\usepackage{fancyvrb}

\pagecolor{black}
\color{white}
\hypersetup{colorlinks=true, linkcolor=cyan, citecolor=cyan, urlcolor=cyan}

\newtheorem{theorem}{Theorem}[section]
\newtheorem{lemma}[theorem]{Lemma}
\newtheorem{corollary}[theorem]{Corollary}
\newtheorem{axiom}[theorem]{Axiom}
\theoremstyle{definition}
\newtheorem{definition}[theorem]{Definition}
\newtheorem{example}[theorem]{Example}
\theoremstyle{remark}
\newtheorem{remark}[theorem]{Remark}

\newcommand{\R}{\mathbb{R}}
\newcommand{\Q}{\mathbb{Q}}
\newcommand{\Z}{\mathbb{Z}}
\newcommand{\N}{\mathbb{N}}
\newcommand{\Rstar}{\mathbb{R}^{\star}}
\newcommand{\eps}{\varepsilon}
\newcommand{\st}{\operatorname{st}}
\newcommand{\halo}{\operatorname{halo}}

\title{Hyperreal Numbers: $\eps \cdot \omega = 1$\\
\large An Algebraic, Computable Introduction to Infinitesimal Calculus}
\author{Notes and a formal (Lean~4) companion development}
\date{\today}

\begin{document}
\maketitle

\begin{abstract}
We give a self-contained, algebraic introduction to the hyperreal numbers
$\Rstar$, aimed at undergraduates who have seen a first calculus course but
not necessarily any mathematical logic. Rather than building $\Rstar$ as an
ultrapower of sequences, we introduce it the way one introduces $\mathbb C
= \R(i)$: as a field extension of $\R$ generated by two symbols, an
infinitesimal $\eps$ and an infinite $\omega$, linked by the single
\emph{gauging law} $\eps\cdot\omega = 1$. From this we derive an algebraic
definition of the derivative, an algebraic Riemann integral, and an
algebraic Dirac delta, and use them to resolve the classical puzzle of
``probability zero but not impossible'' events (the \emph{probabilistic
dart}). Every construction here has an executable, machine-checked
counterpart in a Lean~4 formalization of the same axioms (the
\texttt{HyperList} model), so the paper can double as the mathematical
documentation for that code.
\end{abstract}

\tableofcontents

% =====================================================================
\section{Motivation: a new kind of number}
% =====================================================================

Every student learns to accept the imaginary unit $i$ with $i^2=-1$ as a
useful fiction that turns out to be indispensable — for solving cubics,
for AC circuits, for quantum mechanics. Infinitesimals have had a rockier
history: Newton and Leibniz used them freely, the 19$^{\text{th}}$ century
banished them in favor of $\eps$-$\delta$ limits, and in 1960 Abraham
Robinson showed that infinitesimals can be put on a footing exactly as
rigorous as the reals \cite{robinson}. Nonstandard analysis has since
produced clean, purely algebraic proofs of the core theorems of calculus.

This paper takes the same practical attitude towards infinitesimals that
we take towards $i$: we do not need a specific \emph{construction} to use
$\eps$ productively, we need a specific \emph{fixed number} with clear
algebraic properties. Just as $i$ is not ``an arbitrary square root of
$-1$'' but \emph{the} chosen one, our $\eps$ is not an arbitrary
infinitesimal but a fixed, canonical one, and $\omega := \eps^{-1}$ is its
canonical infinite reciprocal.

\begin{center}
\fbox{\parbox{0.9\textwidth}{\centering
\textbf{The one equation you need to remember:} \quad
$\eps \cdot \omega = 1$
}}
\end{center}

Everything below is a systematic unpacking of what that equation buys us.

% =====================================================================
\section{Axioms for the hyperreal numbers}
% =====================================================================

\subsection{Order axioms}

\begin{axiom}[Canonical infinitesimal]\label{ax:eps}
There is an element $\eps \in \Rstar$ with
\[
0 < \eps < r \qquad \text{for every real } r > 0.
\]
More generally, for every $a \in \R^{+}$, $0 < a\eps < r$ for every real
$r>0$: any positive real multiple of $\eps$ is still smaller than every
positive real.
\end{axiom}

\begin{axiom}[Canonical infinite]\label{ax:omega}
There is an element $\omega \in \Rstar$ with
\[
r < \omega \qquad \text{for every real } r.
\]
\end{axiom}

\begin{axiom}[Gauging]\label{ax:gauge}
\[
\omega := \eps^{-1}, \qquad\text{equivalently}\qquad \eps \cdot \omega = 1.
\]
\end{axiom}

A pleasant consequence: \emph{``infinity plus
one'' is no longer ``infinity''.} Classically one writes $\infty = r\cdot
\infty$ for every $r>0$; here, for $r \ne 1$,
\[
r\cdot\omega \ne \omega,
\]
because $\Rstar$ is an ordered field and $\omega$ is a genuine
'number', not a limit symbol. This is the first hint that hyperreal
arithmetic is \emph{strictly more informative} than the classical
$\pm\infty$.

\subsection{Field extension axioms}

\begin{axiom}\label{ax:field}
$\Rstar$, with $<^{\!*}$ and the functions $+,-,\cdot,{}^{-1}$, is an
ordered field extending $\R$: it satisfies the field axioms, the order
axioms, and every real $r\in \R$ embeds into $\Rstar$ compatibly with
$+,\cdot,<$.
\end{axiom}

\begin{axiom}[Transfer, informal]\label{ax:transfer}
Every first-order statement about $\R$ (built from $+,\cdot,<,0,1$ and
quantifiers over reals) that is true of $\R$ remains true of $\Rstar$ when
the quantifiers range over $\Rstar$ instead — \emph{provided} it does not
secretly quantify over sets of reals (the Archimedean and completeness
properties are excluded on purpose, since $\Rstar$ is neither
Archimedean nor complete).
\end{axiom}

We will not use the Transfer Axiom in its full generality in this paper —
it is what makes nonstandard analysis a complete replacement for
$\eps$-$\delta$ arguments in general, but everything we prove below is
elementary enough to check by direct calculation instead. See
Keisler~\cite{keisler-found} for the precise first-order formulation and
the Elementary Calculus text~\cite{keisler-calc} for the pedagogical case.

\begin{remark}[Existence]
Axioms~\ref{ax:eps}--\ref{ax:transfer} are jointly consistent: an explicit
model is exhibited in Section~\ref{sec:model} as executable Lean~4 code,
and it type-checks. So this is not merely a wish-list — every fact proved
axiomatically below is also a proved \emph{theorem} about a concrete
data structure.
\end{remark}

\subsection{Building the field: $\Rstar = \R(\eps,\omega) = \R(\eps)$}

Because $\omega = \eps^{-1}$, we only need to adjoin $\eps$:
\[
\Rstar \;=\; \R(\eps) \;\cong\; \left\{ \sum_{i} a_i\, \eps^{i} \;:\; a_i \in \R,\ i \in \Z,\ \text{finitely many } a_i \ne 0 \right\},
\]
a field of finite Laurent polynomials in $\eps$ — exactly analogous to
$\mathbb{C} = \R(i) = \{a+bi : a,b\in\R\}$, except the ``basis'' is
infinite: $\{\ldots,\eps^{2},\eps,1,\omega,\omega^{2},\ldots\}$ instead of
$\{1,i\}$. We call the exponent $i$ of a term $a_i\eps^i$ its \textbf{order}:
negative orders are infinitesimal directions, positive orders are infinite
directions, order $0$ is the ordinary real line.

\begin{definition}[Classification]
For $x = \sum_i a_i \eps^i \in \Rstar$, write
\[
x = \underbrace{a_0}_{\text{real part}} \;+\; \underbrace{\textstyle\sum_{i>0} a_i\eps^i}_{\text{infinitesimal part}} \;+\; \underbrace{\textstyle\sum_{i<0} a_i\eps^i}_{\text{infinite part}}.
\]
$x$ is \emph{real} if only $a_0 \ne 0$; \emph{finite} if no negative-index
(infinite) term is present; \emph{infinitesimal} if $a_0=0$ and no
negative-index term is present.
\end{definition}

\begin{definition}[Standard part]
For finite $x$, the \textbf{standard part} $\st(x)$ (also written
$\check x$ or $\operatorname{sh}(x)$) is the real part $a_0$. We extend it
to unbounded $x$ by $\st(x)=+\infty$ if the infinite part is positive,
$\st(x)=-\infty$ if negative.
\end{definition}

\begin{lemma}[$\st$ is a ring homomorphism on finite elements]\label{lem:st-hom}
For finite $x,y \in \Rstar$:
$\st(x+y)=\st(x)+\st(y)$, $\st(x-y)=\st(x)-\st(y)$, $\st(xy)=\st(x)\st(y)$,
and if $\st(y)\ne 0$ then $\st(x/y) = \st(x)/\st(y)$.
\end{lemma}
This is the algebraic engine behind ``taking a limit'': every limit
argument in classical calculus becomes ``compute algebraically in
$\Rstar$, then apply $\st$.''

\begin{definition}[Halo / monad]
$\halo(x) = \{y \in \Rstar : x - y \text{ is infinitesimal or } 0\}$, the
cloud of hyperreals infinitely close to $x$. Two elements with $x-y$
infinitesimal (or $0$) are written $x \approx y$.
\end{definition}

% =====================================================================
\section{The formal model: $\Rstar$ as sorted lists of $(\text{coefficient}, \text{order})$ pairs}\label{sec:model}
% =====================================================================

The axioms above describe \emph{what} $\Rstar$ must satisfy; they do not
by themselves give you something you can compute with. The companion Lean
development builds a concrete, executable field satisfying them.

\subsection{Representation}

An element of \texttt{HyperList} (the reference model, file
\texttt{Hyper/HyperList.lean}) is a finite list of pairs
$(\text{coeff}, \text{order}) \in \Q \times \Q$, representing a finite sum
$\sum (\text{coeff}_k)\,\eps^{\text{order}_k}$, kept in a canonical
\emph{sorted, zero-free, duplicate-free normal form}: for instance
$3 + 2\eps - \omega^2$ is stored as the three pairs $(3,0), (2,-1),
(-1,2)$, sorted by descending order. Addition merges two such lists and
collapses equal-order terms; multiplication is the Cauchy/convolution
product on exponents (orders add, coefficients multiply), matching
ordinary Laurent-polynomial multiplication. Order is \emph{lexicographic
on the leading (highest, i.e.\ ``most infinite'') nonzero term} — this is
exactly what makes $\omega \gg 1 \gg \eps$ provable by comparing leading
terms.

On this concrete type, every fact stated as an axiom in
Section~2 — positivity and gauging of $\eps$ and $\omega$, the additive
identities, the disequalities, and the standard-part function
$\st$ (implemented as ``keep only the order-$0$ term'') — is instead a
proved \emph{theorem}, checked automatically by evaluating both sides of
the decidable statement. In other words, the executable model is not
merely consistent with the axioms; it computes with them.

\subsection{Why not axiomatize directly?}

An earlier experiment declared the hyperreals as an opaque axiomatic
type, exactly mirroring Axioms~\ref{ax:eps}--\ref{ax:transfer}. It is
mathematically the cleanest statement of the theory, but it is \emph{not
executable}: one cannot evaluate $\eps+1$ or numerically test that a
derivative comes out right, because an axiomatized type carries no
algorithm. The list-based model sacrifices some elegance (working with
lists, needing a simplification/normal-form step) in exchange for being
fully computable — which matters for a project whose stated goal is
executable hyperreal algebra, not just a consistency proof.

% =====================================================================
\section{Application 1: the algebraic derivative}
% =====================================================================

Classically, $f'(x) = \lim_{h\to 0} \frac{f(x+h)-f(x)}{h}$. With $\eps$ in
hand we can \emph{skip the limit entirely}:

\begin{definition}[Algebraic derivative]
\[
\partial f(x) \;:=\; \frac{f(x+\eps) - f(x)}{\eps}, \qquad
f'(x) := \st\big(\partial f(x)\big).
\]
\end{definition}

\begin{example}
Let $f(x)=x^2$. Then
\[
\partial f(x) = \frac{(x+\eps)^2 - x^2}{\eps} = \frac{2x\eps + \eps^2}{\eps} = 2x + \eps,
\]
purely by field arithmetic — no limiting process, just algebra in
$\Rstar$. Taking the standard part: $f'(x) = \st(2x+\eps) = 2x$, the
familiar answer, but the intermediate expression $2x+\eps$ carries
\emph{more} information than the classical derivative: it tells you the
exact first-order error of the difference quotient at gauge $\eps$.
\end{example}

\subsection{Derivatives of jump functions: an algebraic Dirac delta}

This is where the hyperreal approach earns its keep — it differentiates
functions that classical calculus can only handle distributionally.

\begin{definition}[Heaviside step]
\[
H(x) := \begin{cases} 1 & x \ge 0 \\ 0 & x < 0\end{cases}.
\]
\end{definition}

Evaluated at the single hyperreal point $x=0$,
\[
\partial H(0) = \frac{H(\eps)-H(0)}{\eps} = \frac{1-1}{\eps} = 0,
\]
which looks wrong until we notice we should instead look at the halo:
approaching from a generic infinitesimal perspective, $H$ jumps from $0$
to $1$ over a step of width $0$ classically, but the correct hyperreal
statement is that the \emph{derivative supported in the halo of $0$}
is $\omega$:
\[
\partial(x \ge 0)\,(\eps) = \omega, \qquad \partial(x\ge0)\,(y) = 0 \text{ for } y \not\approx 0.
\]
Intuitively: the step rises by $1$ over a run of $\eps$, so its slope
\emph{at} the jump is $1/\eps = \omega$ — an honest, finite-in-$\Rstar$
number, not a divergent limit. Define, following the axiomatic sketch in
this project's notes,
\[
\omega_0(x) := \begin{cases}\omega & x \approx 0\\ 0 & \text{else}\end{cases},
\qquad
\boxed{\delta(x) := \tfrac12\,\omega_0(x)}
\]
as the \textbf{algebraic Dirac delta}. One checks it has the two defining
properties of the classical $\delta$-distribution, now as ordinary
Riemann-style integrals (Section~\ref{sec:integral}) instead of formal
symbols:
\[
\int_{-\eps}^{\eps} \delta \;=\; 1, \qquad \delta = \partial H.
\]
Every higher derivative of $H$ is similarly a higher power of $\omega$:
$\partial(x{=}0)(0) = \omega$ (a ``spike''), $\partial^2(x{=}0)(0) =
\omega^2$ (a ``second-order spike''), and so on — dual numbers ($\eps^2=0$)
deliberately \emph{lose} this structure, which is one reason this project
does not adopt the dual-number simplification as its default axiom.

% =====================================================================
\section{Application 2: the algebraic Riemann integral}\label{sec:integral}
% =====================================================================

\begin{definition}[Infinite Riemann sum]
Partition $[a,b]$ into hyperreal-width slices of size $\Delta x$ (any
positive hyperreal, in particular $\Delta x = \eps$), and set
\[
S(f,a,b,\Delta x) \;=\; \sum_{k} f(a+k\Delta x)\,\Delta x
\quad\text{(a finite sum, since $\Rstar$ is a field, not a limit)}.
\]
\end{definition}

\begin{definition}[Algebraic integral]
\[
\int_a^b f(x)\,dx \;:=\; \st\big(S(f,a,b,\eps)\big).
\]
\end{definition}

Because the Riemann sum $S(f,a,b,\Delta x)$ is a genuine \emph{finite
sum}, defined for every positive real $\Delta x$, and $\Rstar$ satisfies
the same first-order properties as $\R$ (Transfer, Axiom~\ref{ax:transfer}),
the same formula is defined for the hyperreal value $\Delta x = \eps$
too — this is not a new definition smuggled in through a limit, it is the
\emph{same} function evaluated at an infinitesimal input, which is
possible in $\Rstar$ precisely because $\Rstar$ is a field extension, not
merely a completion.

\begin{lemma}[Basic properties]\label{lem:integral-basics}
\[
\int_a^b c \;dx = c(b-a), \qquad
\int_a^b f + \int_b^c f = \int_a^c f, \qquad
\frac{d}{dx}\int_a^x f = f(x)\ \text{(FTC)}.
\]
\end{lemma}

\subsection{Gauging the integral: what is $\int \eps$?}

Because $\eps$ itself is a function of the gauge $\Delta x$ (it \emph{is}
the gauge), integrating the constant function $\eps$ against itself gives
back the length of the interval, scaled by the gauging law:
\[
\int_0^\omega \eps \,dx = \omega\cdot\eps - 0\cdot\eps = 1,
\qquad
\int_{-\omega}^{\omega}\eps\,dx = 2.
\]
This ``$\int \eps$'' normalization is the algebraic reason the halo
integral of the algebraic Dirac delta comes out to exactly $1$:
$\int_{-\eps}^{\eps}\omega\,dx = 2$, so $\delta=\tfrac12\omega_0$ integrates
to $1$ as required.

% =====================================================================
\section{Application 3: the probabilistic dart —\\ resolving ``probability zero but not impossible''}
% =====================================================================

\subsection{The classical puzzle}

Throw a dart at a disc of area $A$, uniformly at random. Classical
(Kolmogorov / Lebesgue) probability theory assigns
\[
P(\text{hits the exact point } z) = 0,
\]
for \emph{every} point $z$, even though the dart obviously lands
\emph{somewhere}, and that somewhere is some specific point. The usual
resolution — ``probability $0$ doesn't mean impossible, it means measure
zero'' — is correct but unsatisfying pedagogically: it asks students to
accept that a perfectly certain event (landing on \emph{a} point) is built
out of individually impossible-looking events (landing on \emph{this}
point).

\subsection{The hyperreal resolution}

Replace the real-valued probability measure with a hyperreal-valued one.
In an $n$-dimensional space, the correct atomic unit of probability mass
is $\eps^n$, with the corresponding gauging law $\omega^n\eps^n = 1$
(immediate from Axiom~\ref{ax:gauge}, since $(\omega\eps)^n = \omega^n\eps^n
= 1^n = 1$). Concretely, for a disc of area $A$:

\begin{definition}[Hyperreal dart measures]
\[
\text{pointMass}(A) := \frac{\eps^2}{A}, \qquad
\text{lineMass}(L,A) := \frac{L\cdot\eps}{A} \quad (L = \text{length of the segment}).
\]
\end{definition}

\begin{theorem}[Possible, not impossible]\label{thm:dart-pos}
For $A>0$: $\ \text{pointMass}(A) > 0$. For $A,L>0$:
$\ \text{lineMass}(L,A) > 0$.
\end{theorem}

\begin{theorem}[Classical theory is the shadow]\label{thm:dart-shadow}
For $A\ne 0$: $\ \st(\text{pointMass}(A)) = 0$. For $A,L\ne0$:
$\ \st(\text{lineMass}(L,A)) = 0$.
\end{theorem}

So both events are \emph{genuinely possible} (Theorem~\ref{thm:dart-pos})
\emph{and} classically invisible (Theorem~\ref{thm:dart-shadow}) — the
hyperreal theory strictly refines the classical one rather than
contradicting it: apply $\st$ and you recover exactly Kolmogorov's answer.

\begin{theorem}[Lines are infinitely more likely than points]\label{thm:dart-ratio}
For $A,L>0$:
\[
\text{pointMass}(A) \;<\; \text{lineMass}(L,A), \qquad
\frac{\text{lineMass}(L,A)}{\text{pointMass}(A)} = L\cdot\omega,
\]
an honestly infinite ratio.
\end{theorem}

This matches intuition perfectly: a line is a strictly ``fatter'' target
than a point, and now that fact is a provable \emph{inequality} between
positive numbers, rather than a vacuous $0 < 0$ or an appeal to different
flavors of measure zero. Table~\ref{tab:dims} summarizes the pattern,
which generalizes to any dimension via $\omega^n\eps^n=1$.

\begin{table}[h]
\centering
\begin{tabular}{lllll}
\hline
Space & Object & Atom count & Atom mass & $\st(P)$\\
\hline
1-D & point & $1$ & $\eps$ & $0$ \\
1-D & $[0,1]$ interval & $\omega$ & $\eps$ & $1$ \\
2-D & point & $1$ & $\eps^2$ & $0$ \\
2-D & line, length $L$ & $L\omega$ & $\eps^2$ & $0$ \\
2-D & disc, area $A$ & $A\omega^2$ & $\eps^2$ & $1$ \\
\hline
\end{tabular}
\caption{Dimensional hierarchy of hyperreal dart probabilities. Every row's
total mass is (atom count)$\times$(atom mass); the gauging law
$\omega^n\eps^n=1$ makes the totals for the full space come out to
exactly $1$.}
\label{tab:dims}
\end{table}

\subsection{Formal status}

All four boxed theorems above (\ref{thm:dart-pos}--\ref{thm:dart-ratio})
are proved, with \emph{no} unproven gaps, as concrete theorems about the
executable list-based model described in Section~\ref{sec:model} — not
merely sketched axiomatically. (An earlier, purely axiomatic sketch of the
same statements over an opaque hyperreal type exists as well, but the
version actually machine-checked end to end is the one over the concrete
model, with area and length taken as rational rather than real
parameters, matching that model's scalar field.)

% =====================================================================
\section{Worked exercises}
% =====================================================================

These are pitched at the level of a first course, meant to be solved with
nothing but the field axioms and the gauging law $\eps\omega=1$.

\begin{enumerate}[label=\textbf{Exercise \arabic*.}, leftmargin=2.2cm]
\item Show $\omega^2\eps^2 = 1$ and, more generally, $\omega^n\eps^n=1$ for
every $n\in\N$.
\item Compute $\partial f(x)$ and $f'(x)=\st(\partial f(x))$ for
$f(x)=x^3$. (Answer: $\partial f(x) = 3x^2 + 3x\eps + \eps^2$, so
$f'(x)=3x^2$.)
\item Show that $2\eps$ is infinitesimal and that $\eps+\omega$ is
neither infinitesimal, finite, nor infinite in the usual sense — classify
it correctly using the order decomposition of Section~2.3.
\item Using $\st(x\cdot y) = \st(x)\st(y)$ for finite $x,y$
(Lemma~\ref{lem:st-hom}), show that a positive-real multiple of $\eps$ is
still infinitesimal, and use this to prove $r\cdot\eps < r'$ for all real
$r,r'>0$.
\item (Dart, warm-up) In a 1-D uniform distribution on $[0,\omega]$
(``length $\omega$''), show a single point has mass $\eps$ and the whole
interval has mass $1$. What does this say about $\eps\cdot\omega$
concretely, in probabilistic language?
\item (Dart, main event) Redo Theorem~\ref{thm:dart-ratio} for a disc
versus a full 2-D subregion of area $A' < A$: show
$\text{pointMass}(A) < $ ``area-$A'$-mass'' $= A'\omega^2\eps^2/A$ — wait,
what does that simplify to, and why should you have expected the answer
without computing?
\item Prove $\eps\cdot\eps < \eps$ directly from Axiom~\ref{ax:eps} (hint:
apply the axiom to $r=\eps$ itself, using $a=\eps>0\in\R^+$... but $\eps$
is not real! Explain why the \emph{stated} axiom does not immediately give
this, and instead derive it from $0<\eps<1$ and multiplying both sides of
$\eps<1$ by $\eps>0$).
\end{enumerate}

% =====================================================================
\section{Discussion: what this framework does and does not claim}
% =====================================================================

\begin{itemize}
\item \textbf{It does not claim a new physical theory of probability.}
The hyperreal probability measure is finitely additive on the algebra we
define it on; the delicate countably-additive extension to a full Loeb
measure~\cite{loeb} is a genuinely harder construction, sketched but not
carried out in \texttt{Hyper/HyperProbability.lean} (several \texttt{sorry}s
mark exactly where Loeb measure theory would be needed).
\item \textbf{It does not replace $\eps$-$\delta$ analysis, it
translates it.} Every hyperreal proof here can, in principle, be unwound
into a classical $\eps$-$\delta$ proof via the Transfer Axiom; the payoff
is that the hyperreal proofs are often shorter and more algebraic, and (as
demonstrated by the Lean development) genuinely \emph{executable}.
\item \textbf{The choice $\Rstar=\R(\eps)$ (finite Laurent polynomials) is
a specific, computable model}, not the unique one satisfying
Axioms~\ref{ax:eps}--\ref{ax:transfer}. The classical ultrapower
construction is more general (it supports full transfer for
\emph{all} first-order statements) but is not computable; see
Remark on \S\ref{sec:model}.
\end{itemize}

% =====================================================================
\begin{thebibliography}{9}

\bibitem{robinson} A. Robinson, \emph{Non-standard Analysis}, 1966.

\bibitem{keisler-calc} H. J. Keisler, \emph{Elementary Calculus: An
Infinitesimal Approach}, freely available at
\url{https://people.math.wisc.edu/~hkeisler/calc.html}.

\bibitem{keisler-found} H. J. Keisler, \emph{Foundations of Infinitesimal
Calculus}, \url{https://people.math.wisc.edu/~hkeisler/foundations.pdf}.

\bibitem{loeb} P. A. Loeb, M. P. H. Wolff (eds.), \emph{Nonstandard
Analysis for the Working Mathematician}, Springer, 2015.

\bibitem{goldblatt} R. Goldblatt, \emph{Lectures on the Hyperreals: An
Introduction to Nonstandard Analysis}, Springer GTM 188, 1998.

\end{thebibliography}

\end{document}
